\section{Finite-query doctrine}

\subsection{Mathematical concepts as protocols of finite use}

A weak arithmetic theory cannot conveniently manipulate an actual complete
real line, topological space, scheme, or category as a higher-order object.
That is not necessary for most concrete proofs.  A proof uses such an object
through finitely many observations:

\begin{itemize}
\item an approximation to precision \(p\);
\item values on a finite list of inputs;
\item a finite presentation up to degree or height \(H\);
\item a mesh or simplicial complex at scale \(2^{-p}\);
\item a finite diagram and a requested universal construction;
\item a bounded trace of length \(T\);
\item a finite list of local charts and overlap checks.
\end{itemize}

\begin{definition}[Finite query]
A finite query is a natural-number code that determines every finite demand of
the downstream theorem: precisions, bounds, dimensions, horizons, tolerances,
mesh sizes, input points, and the requested observable.
\end{definition}

The query can be a tuple encoded by pairing or mixed radix.  The important
condition is that it is fixed before the relevant induction or search.

\begin{definition}[Presentation by finite use]
A presented mathematical object is a family of finite certificate problems
indexed by queries, together with:
\begin{enumerate}
\item a primitive-recursive bound on certificate codes;
\item a fixed primitive-recursive verifier;
\item an object-iPRA theorem that accepted certificates answer the query;
\item producer data showing that an accepted certificate exists;
\item coherence or stability theorems relating answers to different queries.
\end{enumerate}
\end{definition}

This definition deliberately separates \emph{meaning} from
\emph{representation}.  Two certificate systems may present the same
observable mathematics if they translate to each other or if their observations
are provably stable.

\subsection{Four output modes}

Not every construction should be forced into a deterministic exact function.
The framework recognizes four modes.

\begin{description}[leftmargin=2.8cm,style=nextline]
\item[Exact deterministic.]
A fixed PR term computes the output and its equations prove the specification.
Examples: integer arithmetic, a fixed matrix determinant, or a bounded sum.

\item[Bounded relational.]
A finite set of codes may be valid.  The least accepted code is selected only
when a concrete representative is needed.  Otherwise the relation itself is
the correct output.

\item[Stable observable.]
Certificates may be unrelated, but every two valid certificates give close or
equal observations.  The observable is well defined without selecting a
coherent representative family.

\item[Conditional.]
The finite verifier is complete only after supplying explicit global data such
as a height bound, analytic constant, separation margin, or covering
representative.  The resulting theorem keeps that data as a premise.
\end{description}

These modes prevent a common constructive error: deriving an exact global
selector from an existence or approximation theorem that justifies only a
relation or a stable quotient.

\subsection{Query-first compilation}

The recommended order is:

\begin{enumerate}
\item state the consumer's finite query;
\item propagate demands backward through the fixed construction;
\item choose numerical bounds for every intermediate certificate;
\item specialize all algorithms and checkers to fixed PR syntax;
\item pack the finite evidence into one code;
\item prove the checker QF and reflected;
\item run bounded search or a direct PR trace;
\item compose object derivations;
\item quantify over the original query only after the QF proof is complete.
\end{enumerate}

This order avoids induction over an ``all precisions at once'' assertion,
prevents hidden dependence on an unbounded evaluator, and makes every missing
producer visible before formalization begins.

\subsection{Efficiency is a refinement, not part of soundness}

Searching all codes below a product or exponential bound can be impractical.
That is acceptable for a normal-form theorem.  Once the specification is fixed,
a specialized polynomial-time or otherwise efficient producer may replace the
brute-force search, provided it proves the same accepted-certificate contract.
The normal form is a proof-theoretic baseline and interoperability layer, not a
claim about optimal algorithms.
